\subsection{Contraste de la constante de Planck reducida}
\label{sec:revision-planck-contraste}

El Sistema Internacional vigente desde el 20 de mayo de 2019 fija
exactamente \(h=6.62607015\times10^{-34}\,\mathrm{J\,s}\)
\cite{bipm2018}. Por tanto, su constante reducida es
\[
 \hbar_{\mathrm{SI}}=\frac{6.62607015}{2\pi}\,10^{-34}\,
 \mathrm{J\,s}
 =1.0545718176461563912\ldots\times10^{-34}\,\mathrm{J\,s}.
\]
Los puntos suspensivos corresponden a la expansión de una expresión
exacta, no a una incertidumbre experimental de \(h\) en el SI actual.
Este valor de referencia se utiliza aquí después de la evaluación de
los lectores de HMT; su función es comparativa.

Las dos coordenadas del retorno del artículo son
\begin{align*}
 H_5&=1.0545718176461544696\ldots,\\
 \eta_{\mathrm{ret}}&=1.0545718169150390611\ldots.
\end{align*}
En la identificación de \(\mathcal U_S\) con \(\mathrm{J\,s}\), la
diferencia relativa de la sección anterior respecto de
\(\hbar_{\mathrm{SI}}\) es aproximadamente \(-1.82\times10^{-15}\).
La de la sección posterior es aproximadamente \(-6.93\times10^{-10}\).
La segunda diferencia contiene el efecto del retorno que utiliza el
corrector electrónico; no equivale al error de evaluación de la
primera sección.

La precisión de estas comparaciones permite reconocer la proximidad
cuantitativa de la construcción de \(H_5\) y, simultáneamente, distinguir
proximidad de identidad exacta. En el presente manuscrito las expresiones
exhibidas no proporcionan una igualdad exacta con \(h/(2\pi)\).
El resultado demostrado sobre \(\Sigma_H\) es su década \(-34\);
el resultado analítico sobre \(H_5\) es la evaluación del carácter
declarado. La identificación física se examina en esta coordenatización y con las
diferencias explícitas que acabamos de dar. El valor SI no interviene
en la definición del lector de norma y autoescala ni en la recurrencia regional.
